\section{Nonuniform calibration and metastable occupation}
\label{sec:occupation-stability}
A small uniform defect at every possible history is unnecessary. What the converse uses is the predictable defect accumulated before actual commitment, including rejected productive calls. The following criterion is stated at raw call boundaries and covers adaptive policies.

\begin{proposition}[A drift certificate for accumulated defects]\label{prop:potential}
Let $d_i\ge0$ be a predictable defect charge at call $i$, zero when that call is not productive. Suppose there are nonnegative integrable potentials $V_i$ on a filtration containing the physical and observer history, with $V_0$ integrable, and a constant $\lambda\ge0$ such that for every permitted current command, on $\{i\le\tau\}$,
\begin{equation}\label{eq:potential-drift}
 d_i+\E_*[V_i\one_{\{i<\tau\}}\mid\mathcal F_{i-1},a_i]
 \le V_{i-1}+\lambda\E_*[c_i\mid\mathcal F_{i-1},a_i].
\end{equation}
After termination set the potential to zero. Then every legal procedure with $\E_*C_\tau\le N$ has
\begin{equation}\label{eq:potential-bound}
 \E_*\sum_{i\le\tau}d_i\le \E_*V_0+\lambda N.
\end{equation}
The potentials are analytic physical-state certificates. They are not additional memory available to the observer. The same assertion holds if the conditional expectation of a nonnegative physical defect dominates the visible-history defect; apply the bound to the physical defect and use the tower property.
\end{proposition}
\begin{proof}
For a deterministic $k$, multiply \eqref{eq:potential-drift} by $\one_{\{i\le\tau\}}$, integrate, and sum from $1$ to $k$. Conditional measurability of the decision to make the call makes all interior potential terms cancel. The remaining terminal term is $\E_*[V_k\one_{\{k<\tau\}}]\ge0$. Drop it, bound the partial expected cost by $N$, and let $k\to\infty$ by monotone convergence. If commands are randomized, integrate the commandwise inequality first. The conditional-dominance version follows by applying the tower property to each nonnegative summand before summing. No optional stopping theorem with unverified uniform integrability is needed.
\end{proof}

\begin{theorem}[Metastable calibration with explicit rare-mode mass]\label{thm:metastable}
Suppose the actual reference contains a hidden mode $D_i\in\{0,1\}$ at raw call boundaries. For every actual history and legal command, its one-call transition obeys
\[
 \Pp_*(D_i=1\mid D_{i-1}=0,\mathcal F_{i-1},a_i)\le u,
 \qquad
 \Pp_*(D_i=0\mid D_{i-1}=1,\mathcal F_{i-1},a_i)\ge v>0.
\]
Here $0\le u\le1$, $0<v\le1$, and $\pi_0=\Pp_*(D_0=1)$ is the actual preparation mass. Every raw call costs at least $c_*>0$. Conditional productive mass and task defects are bounded by the visible conditional expectations of
$a_g+a_b\one_{\{D_{i-1}=1\}}$ and
$e_g+e_b\one_{\{D_{i-1}=1\}}$, respectively. Let $b_-=N/\kappa_-$. Then Theorem~\ref{thm:reference-stability} applies with
\begin{align}
 \mathfrak a_N&=a_gb_-+\frac{a_b}{v}\pi_0+
                         \frac{a_bu}{vc_*}N,\label{eq:metastable-mass}\\
 \mathfrak e_N&=e_gb_-+\frac{e_b}{v}\pi_0+
                         \frac{e_bu}{vc_*}N.\label{eq:metastable-task}
\end{align}
All competing adaptive stopping rules are covered. The bad-mode defects need not be small, and $v$ may be small; their residence time and preparation mass remain explicitly charged. The matching geometric conclusion of Corollary~\ref{cor:reference-geometry} follows whenever these bounds satisfy its inequalities.
\end{theorem}
\begin{proof}
For the bad-mode mass charge use $V(D)=(a_b/v)\one_{\{D=1\}}$. In mode zero its conditional increase is at most $a_bu/v$. In mode one its conditional decrease is at least $a_b$, canceling the bad-mode charge. Hence
\[
 a_b\one_{\{D_{i-1}=1\}}+\E_*[V(D_i)-V(D_{i-1})\mid\mathcal F_{i-1},a_i]
 \le a_bu/v.
\]
Killing the potential on commitment can only decrease the left side. Since $c_i\ge c_*$, Proposition~\ref{prop:potential} bounds the sum over all calls by $a_b\pi_0/v+a_buN/(vc_*)$. Charging only productive calls is smaller. The good baseline contributes at most $a_g\E_*A_\tau\le a_gb_-$. This proves \eqref{eq:metastable-mass}; replace $a$ by $e$ for \eqref{eq:metastable-task}. Observations may reveal or bias the mode on surviving histories; conditioning and the drift estimate already include that selection. Neither stationary preparation nor independence of the stopping time from the mode is assumed.
\end{proof}

This is a substantive nonuniform class. For example take $a_g=e_g=0$, a fixed bad-mode defect, $u=\eta v$, and initial mass $\pi_0=\eta$, with $\eta,v>0$. The rowwise worst defect remains fixed as $\eta\to0$, while the actual accumulated bound is of order $\eta/v+\eta N/c_*$. Long bad episodes are permitted. The transient term cannot be silently removed by replacing actual preparation with stationary mass. Such a mode can modulate any positive report--audit instrument, including an atomic or changing-rank instrument; its mixture weights in this proof are the actual mode probabilities.

\begin{proposition}[Selection amplifies reference error in both directions]\label{prop:reference-sharpness}
For every $b\ge1$ and $0\le e\le1/b$ there are positive prepared comparison and actual productive kernels with identical report laws, one unit of cost per call, and uniform task defect $e$, for which
\[
 \Qm_m(b)=be,\qquad R_\cp(b,m)=R_\evt(b,m)=0
\]
for all $m\ge1$. The comparison need not belong to the actual singleton family. Reversing the two kernels gives $\Qm_m(b)=0$ and actual minimax risk $be$. Thus a term of order $be$ is unavoidable on both sides of a general transfer statement, even with unrestricted terminal width. The appropriate state-stopped constructions use one additional live label.
\end{proposition}
\begin{proof}
Let each call draw an independent uniform report $x\in[0,1]$, and let $D$ have mass $1/b$. Audits are unobserved, conditionally mean zero and independent of prior history. Outside $D$, in both experiments the audit is a fresh fair sign. On $D$ put a fresh fair sign of magnitude $\sqrt{be}$ in the comparison and put zero in the actual experiment. Both first moments vanish. For every $q\in[0,1]$ and every report-dependent finite readout, the difference of unnormalized squared losses is
$be\int_D q\,d\nu\le e$.
The comparison's residual variance is everywhere at least $be$, and selecting $D$ with output zero attains that variance at expected cost $b$. Its trimmed profile is therefore $be$ for every positive width. In the actual experiment the same rule has risk zero. One event-stopped state can reject by remaining at itself and commit by emitting a final event while retaining that same state; two intrinsically halt-typed states are required and sufficient. On reversing the kernels, every actual conditional residual variance is at least $be$, so every almost-surely committing strategy has risk at least $be$, and the same rule attains it. The audit amplitude, not its report marginal, has changed. This is an all-program argument, not a failure of one proposed sampler.
\end{proof}
